\chapter[Языковые модели и постобучение]{Языковые модели и постобучение}
\label{ch:lm-post-training}

\begin{quote}
\itshape
Горький урок шестидесяти лет исследований в области ИИ состоит в том, что общие методы, использующие вычисления, в конечном счёте наиболее эффективны. Языковые модели --- чистейшее выражение этого урока из всех открытых на сегодня.
\end{quote}

\bigskip

Предшествующие десять глав построили аппарат обучения с подкреплением с первых принципов: марковские процессы принятия решений, динамическое программирование, методы Монте-Карло и разностного временного обучения, аппроксимация функций, градиенты стратегии и алгоритмы актор-критик. Глава~\ref{ch:reward-models} расширила этот аппарат на задачу обучения сигнала вознаграждения по данным предпочтений людей. Всё это в каком-то смысле было прелюдией. Последующие главы\,---\,о проксимальной оптимизации стратегии для языковых моделей, о прямой оптимизации предпочтений, о групповой относительной оптимизации стратегии\,---\,применяют RL к одному из наиболее значимых вычислительных артефактов нашей эпохи: большой языковой модели.

Прежде чем эти применения можно развить точно, читателю необходимо рабочее понимание того, что такое языковая модель, как она строится и почему RL --- правильный инструмент для финальных стадий её обучения. Эта глава предоставляет такой фундамент. Она не является исчерпывающим обзором обработки естественного языка\,---\,для этого существует несколько отличных учебников,\,---\,но скорее самодостаточным мостиком между RL-теорией Части~II и методами RLHF Части~III.

Начнём с архитектуры Трансформера (Раздел~\ref{sec:transformer}), нейронной основы всех современных языковых моделей, затем формализуем авторегрессионное языковое моделирование как статистическую задачу (Раздел~\ref{sec:autoregressive-lm}). Раздел~\ref{sec:lm-mdp} устанавливает ключевую концептуальную связь: генерация языка --- это марковский процесс принятия решений, а языковая модель сама является стратегией. Вооружившись этим соответствием, Разделы~\ref{sec:pretraining}--\ref{sec:post-training-pipeline} описывают трёхэтапный рецепт обучения, используемый на практике: массивное предобучение, дообучение с учителем (SFT) и выравнивание с помощью обучения с подкреплением. Раздел~\ref{sec:synthetic-data} рассматривает растущую роль синтетических данных, ставших незаменимыми на каждом этапе постобучения. Раздел~\ref{sec:lm-eval} охватывает методологию оценки, а Раздел~\ref{sec:lm-summary} показывает место этой главы в общей логике книги.

% ===========================================================
\section{Архитектура Трансформера}
\label{sec:transformer}
\index{Transformer architecture}

Современные большие языковые модели почти повсеместно реализованы на основе архитектуры Трансформера, предложенной Vaswani et al.\ (2017). Для последующего RL-анализа достаточно понимать Трансформер на концептуальном уровне\,---\,без погружения во все детали учебника по NLP. Мы опишем четыре его ключевых компонента: само-внимание, позиционное кодирование, блоки прямой связи и нормализацию слоёв.

\subsection{Само-внимание}
\index{self-attention}

Центральной инновацией Трансформера является механизм \textbf{само-внимания}, позволяющий каждой позиции в последовательности непосредственно обращать внимание на каждую другую. Для последовательности из $n$ представлений токенов $\mathbf{x}_1, \ldots, \mathbf{x}_n \in \R^d$, упакованных в матрицу $X \in \R^{n \times d}$, масштабированное скалярное произведение вычисляет

\begin{equation}
\label{eq:self-attention}
\operatorname{Attention}(Q, K, V) = \operatorname{softmax}\!\left(\frac{QK^\top}{\sqrt{d_k}}\right) V,
\end{equation}

где $Q = XW_Q$, $K = XW_K$, $V = XW_V$ --- проекции запроса, ключа и значения с обученными весовыми матрицами $W_Q, W_K, W_V \in \R^{d \times d_k}$. Деление на $\sqrt{d_k}$ предотвращает рост скалярных произведений настолько, что softmax насыщается.

Интуитивно, каждый вектор запроса $\mathbf{q}_i$ сравнивается со всеми векторами ключей $\mathbf{k}_j$ для получения оценки внимания, которая затем используется для формирования взвешенного среднего векторов значений. Токен, представляющий слово «он», может сильно обращать внимание на существительное, которое оно заменяет, независимо от расстояния между ними. Это неограниченное рецептивное поле\,---\,в отличие от локальных окон свёрточных сетей или последовательного узкого места рекуррентных сетей\,---\,делает Трансформер столь мощным для языка.

На практике $h$ голов внимания вычисляются параллельно и конкатенируются:
\[
\operatorname{MultiHead}(X) = \operatorname{Concat}(\mathrm{head}_1, \ldots, \mathrm{head}_h)\,W_O,
\]
где каждая $\mathrm{head}_i = \operatorname{Attention}(XW_{Q,i}, XW_{K,i}, XW_{V,i})$ и $W_O \in \R^{hd_k \times d}$. Множество голов позволяет модели одновременно обращать внимание на разные аспекты контекста: одна голова может улавливать синтаксическое согласование, другая --- кореференцию, третья --- семантическое сходство.

\subsection{Позиционное кодирование}
\index{positional encoding}

Само-внимание \emph{эквивариантно к перестановкам}: если позиции токенов перемешаны, оценки внимания перемешиваются аналогично, и модель не может отличить «кошка сидела на коврике» от какой-либо его перестановки. Позиционная информация должна быть введена явно.

В исходном Трансформере использовались фиксированные синусоидальные кодировки, добавляемые к эмбеддингам токенов. Современные модели, как правило, предпочитают \textbf{ротационные позиционные эмбеддинги} (RoPE; Su et al., 2021) или \textbf{ALiBi} (Press et al., 2022), кодирующие относительные, а не абсолютные позиции и лучше обобщающиеся на длины последовательностей, не встречавшиеся при обучении. Для наших целей достаточно знать, что после позиционного кодирования модель имеет доступ к порядку токенов.

\subsection{Блоки прямой связи и нормализация слоёв}
\index{feed-forward block}\index{layer normalisation}

Каждый \emph{слой} Трансформера состоит из подслоя само-внимания, за которым следует позиционно-точечный подслой прямой связи. Блок прямой связи применяет два линейных преобразования с нелинейностью между ними:
\[
\operatorname{FFN}(\mathbf{x}) = W_2\,\sigma(W_1 \mathbf{x} + \mathbf{b}_1) + \mathbf{b}_2,
\]
где $\sigma$ --- как правило, функция активации GELU. Внутреннее измерение обычно в четыре раза превышает измерение модели: $W_1 \in \R^{4d \times d}$, $W_2 \in \R^{d \times 4d}$.

\textbf{Нормализация слоёв} (LayerNorm; Ba et al., 2016) применяется перед каждым подслоем (вариант \emph{pre-norm}, используемый в большинстве современных моделей):
\[
\operatorname{LayerNorm}(\mathbf{x}) = \frac{\mathbf{x} - \mu}{\sigma + \epsilon} \odot \boldsymbol{\gamma} + \boldsymbol{\beta},
\]
где $\mu$ и $\sigma$ --- среднее и стандартное отклонение, вычисленные по измерению признаков, а $\boldsymbol{\gamma}, \boldsymbol{\beta}$ --- обученные параметры масштабирования и сдвига. LayerNorm стабилизирует обучение, предотвращая бесконтрольный рост активаций по слоям.

Остаточное соединение охватывает каждый подслой, так что полное вычисление в одном слое Трансформера таково:
\begin{align*}
\mathbf{h}^{(l)} &= \mathbf{h}^{(l-1)} + \operatorname{MultiHead}(\operatorname{LayerNorm}(\mathbf{h}^{(l-1)})),\\
\mathbf{h}^{(l+1)} &= \mathbf{h}^{(l)} + \operatorname{FFN}(\operatorname{LayerNorm}(\mathbf{h}^{(l)})).
\end{align*}
Стек из $L$ таких слоёв и проекция представления верхнего слоя на словарь задают выходное распределение модели. Рисунок~\ref{fig:transformer-block} иллюстрирует один такой слой.

\begin{figure}[ht]
\centering
% Чистый вертикальный слой Трансформера: основной поток слева, остаточные переходы справа
\begin{tikzpicture}[
  node distance=0.55cm,
  box/.style={draw, rounded corners=3pt, minimum width=3.4cm, minimum height=0.65cm,
              fill=blue!8, font=\small, align=center},
  add/.style={draw, circle, minimum size=0.55cm, fill=orange!20, font=\small,
              inner sep=0pt},
  arr/.style={-{Stealth[length=5pt]}, thick},
  lbl/.style={font=\footnotesize\itshape}
]

% ---- Основной вертикальный поток ----
\node[lbl]                       (in)  {$\mathbf{h}^{(l-1)}$};
\node[box,  below=0.45cm of in]  (ln1) {LayerNorm};
\node[box,  below=0.55cm of ln1] (mha) {Многоголовое само-внимание};
\node[add,  below=0.55cm of mha] (a1)  {$+$};
\node[box,  below=0.55cm of a1]  (ln2) {LayerNorm};
\node[box,  below=0.55cm of ln2] (ffn) {Прямая связь (FFN)};
\node[add,  below=0.55cm of ffn] (a2)  {$+$};
\node[lbl,  below=0.45cm of a2]  (out) {$\mathbf{h}^{(l)}$};

% ---- Стрелки основного потока ----
\draw[arr] (in)  -- (ln1);
\draw[arr] (ln1) -- (mha);
\draw[arr] (mha) -- (a1);
\draw[arr] (a1)  -- (ln2);
\draw[arr] (ln2) -- (ffn);
\draw[arr] (ffn) -- (a2);
\draw[arr] (a2)  -- (out);

% ---- Остаточный переход 1: вход -> a1 (минуя LayerNorm + MHA) ----
\draw[arr, rounded corners=4pt]
  (in.east)   -- ++(2.1,0)
              -- ++(0,-3.33)
              -- (a1.east);

% ---- Остаточный переход 2: выход a1 -> a2 (минуя LayerNorm + FFN) ----
\draw[arr, rounded corners=4pt]
  (a1.east)   -- ++(2.1,0)
              -- ++(0,-3.50)
              -- (a2.east);

% ---- Метки переходов ----
\node[lbl, above right=0.02mm and 7.5mm of a1] {skip};
\node[lbl, below right=0.02mm and 7.5mm of a2] {skip};

\end{tikzpicture}
\caption{Один слой декодера Трансформера. Входное представление $\mathbf{h}^{(l-1)}$ проходит через нормализацию слоёв перед многоголовым само-вниманием и вновь перед блоком прямой связи. Остаточные соединения (узлы $+$) позволяют градиентам течь непосредственно от выхода к входу, обеспечивая устойчивое обучение глубоких сетей.}
\label{fig:transformer-block}
\end{figure}

\begin{remark}
Для языкового моделирования Трансформер использует архитектуру \emph{только декодера} с \textbf{причинной маской}\index{causal masking}: во время само-внимания позиции $i$ запрещено обращать внимание на позиции $j > i$. Это гарантирует, что предсказание токена $t$ зависит только от токенов $1, \ldots, t-1$, сохраняя авторегрессионную структуру. Маска реализуется установкой $QK^\top_{ij} = -\infty$ при $j > i$ перед применением softmax, что обнуляет соответствующие веса внимания.
\end{remark}

Масштаб современных Трансформеров поражает: GPT-3 имеет 175 миллиардов параметров в 96 слоях; LLaMA~3 70B имеет 80 слоёв со скрытым измерением 8192. Несмотря на масштаб, все эти модели разделяют одни и те же четыре компонента, описанных выше.

% ===========================================================
\section{Авторегрессионное языковое моделирование}
\label{sec:autoregressive-lm}
\index{autoregressive language modelling}

\textbf{Языковая модель} --- это распределение вероятностей на последовательностях токенов. Обозначим \textbf{словарь}\index{vocabulary} как $\mathcal{V}$ --- конечное множество из $|\mathcal{V}|$ токенов, обычно от 32~000 до 128~000 элементов после байт-парного кодирования (BPE; Sennrich et al., 2016). Последовательность токенов $\mathbf{w} = (w_1, w_2, \ldots, w_T)$ с $w_t \in \mathcal{V}$ может представлять предложение, документ, диалог или фрагмент кода.

\begin{definition}[Авторегрессионная языковая модель]
\label{def:autoregressive-lm}
\index{autoregressive language model}
\textbf{Авторегрессионная языковая модель} с параметрами $\theta$ --- это параметризованное распределение на последовательностях токенов, определённое по правилу цепочки вероятностей:
\begin{equation}
\label{eq:lm-chain-rule}
p_\theta(w_1, \ldots, w_T) = \prod_{t=1}^{T} p_\theta(w_t \mid w_1, \ldots, w_{t-1}),
\end{equation}
где каждый множитель $p_\theta(w_t \mid w_{1:t-1})$ --- распределение на $\mathcal{V}$, вычисленное из контекста $(w_1, \ldots, w_{t-1})$.
\end{definition}

Факторизация~\eqref{eq:lm-chain-rule} точна по правилу цепочки и не вносит приближений. Трансформер реализует каждый условный $p_\theta(w_t \mid w_{1:t-1})$, отображая контекст в вектор \textbf{логитов}\index{logits} $\mathbf{z}_t \in \R^{|\mathcal{V}|}$ и применяя softmax:
\begin{equation}
\label{eq:lm-softmax}
p_\theta(w_t \mid w_{1:t-1}) = \operatorname{softmax}\!\left(\mathbf{z}_t / \tau\right),
\end{equation}
где $\tau > 0$ --- параметр \textbf{температуры}\index{temperature}. При $\tau = 1$ (по умолчанию) модель сэмплирует пропорционально вероятностям softmax. При $\tau \to 0$ распределение концентрируется на единственном токене с наибольшей вероятностью (жадное декодирование). Температуры выше 1 выравнивают распределение, увеличивая разнообразие за счёт связности.

\subsection{Цель обучения: кросс-энтропийная потеря}
\index{cross-entropy loss!language modelling}

Имея корпус из $N$ документов, где документ $i$ содержит токены $(w_1^{(i)}, \ldots, w_{T_i}^{(i)})$, цель обучения --- \textbf{предсказание следующего токена}: минимизировать среднее отрицательное логарифмическое правдоподобие

\begin{equation}
\label{eq:ntp-loss}
\mathcal{L}_{\mathrm{NTP}}(\theta) = -\frac{1}{\sum_i T_i}\sum_{i=1}^{N}\sum_{t=1}^{T_i} \log p_\theta\!\left(w_t^{(i)} \mid w_{1:t-1}^{(i)}\right).
\end{equation}

Это кросс-энтропия между эмпирическим распределением данных и предсказанным распределением модели. Поскольку истинный следующий токен всегда является унитарным вектором над $\mathcal{V}$, кросс-энтропия сводится к отрицательной логарифмической вероятности верного токена. Эквивалентно, минимизация~\eqref{eq:ntp-loss} максимизирует совместную вероятность, присваиваемую моделью обучающему корпусу.

\begin{definition}[Перплексия]
\label{def:perplexity}
\index{perplexity}
\textbf{Перплексия} языковой модели $p_\theta$ на тестовой последовательности $(w_1, \ldots, w_T)$ равна
\begin{equation}
\label{eq:perplexity}
\mathrm{PPL}(\theta) = \exp\!\left(-\frac{1}{T}\sum_{t=1}^{T}\log p_\theta(w_t \mid w_{1:t-1})\right) = \exp(\mathcal{L}_{\mathrm{NTP}}).
\end{equation}
Перплексия --- это геометрическое среднее обратных вероятностей, присваиваемых моделью каждому токену. Перплексия~$k$ означает, что модель в среднем так же неопределённа, как если бы ей приходилось выбирать равномерно из $k$ равновероятных токенов. Современные большие языковые модели достигают перплексий ниже 10 на стандартных бенчмарках по сравнению с базовым значением случайного ответа $|\mathcal{V}| \approx 50~000$.
\end{definition}

\subsection{Законы масштабирования}
\index{scaling laws}

Одним из наиболее практически важных эмпирических открытий в языковом моделировании является то, что качество (измеренное потерей на валидации) следует предсказуемым \textbf{законам масштабирования} в зависимости от размера модели $N$, размера датасета $D$ и вычислительного бюджета $C \approx 6ND$. Kaplan et al.\ (2020) установили, что потеря на валидации убывает как степенная функция от каждого из этих количеств по отдельности. Hoffmann et al.\ (2022)\,---\,в том, что сегодня известно как \textbf{законы масштабирования Chinchilla}\index{Chinchilla scaling laws}\,---\,показали, что предыдущие крупные модели (включая GPT-3) были значительно недообучены относительно своего размера: оптимальное по вычислениям обучение требует примерно 20 токенов на параметр модели. Результат Chinchilla означает, что модель с $N$ параметрами следует обучать приблизительно на $20N$ токенах для минимизации потерь при данном вычислительном бюджете. Это наблюдение изменило подход к проектированию передовых моделей: LLaMA~3 8B, например, обучалась на 15 триллионах токенов\,---\,значительно больше оптимального по Chinchilla количества, что отражает практическую важность эффективности вывода над эффективностью обучения.

% ===========================================================
\section{Генерация языка как марковский процесс принятия решений}
\label{sec:lm-mdp}
\index{language model!as MDP}

Описанный выше авторегрессионный процесс генерации имеет элегантную интерпретацию как марковский процесс принятия решений. Это соответствие является концептуальным фундаментом каждого алгоритма RLHF в Части~III: как только мы распознаём языковую модель как стратегию, весь RL-аппарат, разработанный в Главах~\ref{ch:pg} и~\ref{ch:actor-critic}, применяется непосредственно.

\begin{definition}[MDP уровня токенов]
\label{def:token-mdp}
\index{token-level MDP}
Пусть $x = (x_1, \ldots, x_m)$ обозначает \textbf{промпт}\index{prompt} (последовательность из $m$ токенов), а $y = (y_1, \ldots, y_T)$ --- \textbf{ответ}\index{response}, подлежащий генерации, с $y_t \in \mathcal{V}$. \textbf{MDP уровня токенов} --- это кортеж $\mathcal{M} = (\cS, \cA, P, R, \gamma)$, где:
\begin{itemize}[leftmargin=*]
  \item \textbf{Пространство состояний} $\cS$: состояние на шаге $t$ --- полный контекст, увиденный к этому моменту,
  \[
  s_t = (x_1, \ldots, x_m, y_1, \ldots, y_{t-1}),
  \]
  то есть конкатенация промпта и всех токенов, сгенерированных до (но не включая) шага $t$. Начальное состояние --- $s_0 = x$.
  \item \textbf{Пространство действий} $\cA = \mathcal{V}$: действие на шаге $t$ --- следующий токен $a_t = y_t$. Пространство действий имеет мощность $|\mathcal{V}| \sim 10^4$--$10^5$.
  \item \textbf{Динамика переходов} $P$: детерминированная,
  \[
  s_{t+1} = (x, y_1, \ldots, y_t) = s_t \concat y_t,
  \]
  где $\concat$ обозначает конкатенацию. Переход полностью определён текущим состоянием и действием.
  \item \textbf{Функция вознаграждения} $R$: разреженная и терминальная,
  \[
  R(s_t, a_t) = \begin{cases} r(x, y) & \mbox{если } a_t = \textsc{eos} \mbox{ или } t = T, \\ 0 & \mbox{иначе,} \end{cases}
  \]
  где $r : \mathcal{X} \times \mathcal{Y} \to \R$ --- скалярная оценка полного ответа (например, обученная модель вознаграждения или балл верифицируемой корректности), а \textsc{eos} --- токен конца последовательности.
  \item \textbf{Коэффициент дисконтирования} $\gamma = 1$ (эпизодическая задача с единственным терминальным вознаграждением).
  \item \textbf{Стратегия} $\pi_\theta$: сама языковая модель,
  \[
  \pi_\theta(a_t \mid s_t) = p_\theta(y_t \mid x, y_{1:t-1}).
  \]
\end{itemize}
\end{definition}

\begin{proposition}
\label{prop:lm-return}
Суммарная доходность в MDP уровня токенов с $\gamma = 1$ равна вознаграждению, присвоенному полному ответу:
\begin{equation}
\label{eq:lm-return}
G_0 = \sum_{t=0}^{T} R(s_t, a_t) = r(x, y).
\end{equation}
Следовательно, цель RLHF
\begin{equation}
\label{eq:rlhf-objective}
J(\theta) = \E_{x \sim \mathcal{D},\; y \sim \pi_\theta(\cdot \mid x)}\bigl[r(x, y)\bigr]
\end{equation}
--- это в точности ожидаемая доходность стратегии $\pi_\theta$ в MDP уровня токенов.
\end{proposition}

\begin{proof}
Все промежуточные вознаграждения по определению нулевые; единственное ненулевое слагаемое в сумме --- $R(s_T, a_T) = r(x, y)$. Взятие ожидания по $x \sim \mathcal{D}$ и распределению траектории, порождённому $\pi_\theta$, непосредственно даёт~\eqref{eq:rlhf-objective} из определения ожидаемой доходности.
\end{proof}

Эта MDP-формулировка обнаруживает три структурные особенности, отличающие генерацию языка от классических RL-задач Части~II.

\begin{remark}[Отличительные черты MDP уровня токенов]
\label{rem:lm-mdp-features}
\begin{enumerate}[leftmargin=*]
  \item \textbf{Огромное пространство действий.} При $|\cA| \sim 10^5$ хранить функцию ценности $Q(s, a)$ явно для всех действий вычислительно нецелесообразно. Это исключает табличные методы (Глава~\ref{ch:dp}) и большинство вариантов Q-обучения, мотивируя подходы на основе градиента стратегии и актор-критик из Глав~\ref{ch:pg}--\ref{ch:actor-critic}.
  \item \textbf{Длинный горизонт с разреженным вознаграждением.} Типичный ответ содержит 200--2000 токенов. Единственное терминальное вознаграждение должно быть приписано сотням предшествующих действий\,---\,проблема присвоения кредита в её крайней форме. Модели вознаграждения по процессу (Раздел~\ref{sec:prm} в Главе~\ref{ch:reward-models}) решают её, предоставляя обратную связь на уровне шагов.
  \item \textbf{Детерминированные переходы.} После добавления токена состояние определяется с полной определённостью. Всё стохастическое происходит от самой стратегии, а не от среды. Это упрощает некоторые теоретические анализы, но затрудняет исследование: модель должна полагаться на своё собственное распределение сэмплирования для посещения новых завершений.
\end{enumerate}
\end{remark}

\begin{example}[Подсчёт токенов в MDP]
\label{ex:token-mdp-example}
Рассмотрим промпт $x = $ «\texttt{Какова столица Франции?}» (6 токенов) и ответ $y = $ «\texttt{Столицей Франции является Париж.}» (7 токенов). Эпизод имеет $T = 7$ шагов. На шаге $t = 5$ состояние --- $s_5 = (x, \allowbreak \texttt{Столицей}, \allowbreak \texttt{Франции}, \allowbreak \texttt{является}, \allowbreak \texttt{Париж})$. Действие --- $a_5 = \texttt{.}$. Вознаграждение нулевое для всех $t < 7$, и $r(x, y)$ даётся при $t = 7$ моделью вознаграждения, оценивающей полный ответ. Если модель вознаграждения оценивает фактическую корректность и присваивает $r = +1.4$, то $G_0 = 1.4$ независимо от того, какие промежуточные токены были сгенерированы.
\end{example}

% ===========================================================
\section{Предобучение}
\label{sec:pretraining}
\index{pre-training}

Первый этап построения способной языковой модели --- \textbf{предобучение}: обучение случайно инициализированного Трансформера на огромном корпусе текста с использованием цели предсказания следующего токена~\eqref{eq:ntp-loss}. Предобучение --- безусловно наиболее вычислительно затратный этап и обычно потребляет более 95\% общего бюджета обучения передовой модели.

\subsection{Данные}

Корпусы для предобучения собираются из разнообразных источников в интернете и институциональных коллекций. Общие компоненты включают:

\begin{description}[leftmargin=!, labelwidth=3.5cm, align=parleft]
  \item[Common Crawl] Ежемесячные снимки веба; после агрессивного дедублирования и фильтрации качества обычно составляет 50--70\% обучающих токенов.
  \item[Книги] Project Gutenberg, Books3 и аналогичные коллекции; богаты длинным структурированным текстом со связным повествованием.
  \item[Научная литература] ArXiv, PubMed, Semantic Scholar; критически важна для способностей к техническому рассуждению.
  \item[Код] Репозитории GitHub на десятках языков программирования; сильно коррелирует с улучшенным математическим и структурированным рассуждением.
  \item[Википедия и Викиданные] Высококачественный энциклопедический текст с фактическими утверждениями; как правило, избыточно представлен относительно его исходного размера.
\end{description}

Общее число токенов для современных передовых моделей составляет от одного до пятнадцати триллионов. LLaMA~3 обучалась на 15 триллионах токенов; обучающие данные GPT-4 не раскрываются публично, но, по имеющимся сведениям, сопоставимы по масштабу. Сборка такого корпуса требует значительных инженерных усилий: удаление почти-дубликатов (с использованием хеширования MinHash с локальной чувствительностью), идентификация языка, фильтрация качества (отклонение документов ниже порога перплексии, установленного небольшой эталонной моделью) и очистка персональных данных.

\subsection{Потеря предобучения и эмерджентные способности}

Предобучение оптимизирует~\eqref{eq:ntp-loss} с использованием вариантов AdamW (Loshchilov \& Hutter, 2019) с косинусным расписанием скорости обучения. Градиентные контрольные точки, арифметика смешанной точности (bfloat16) и распределённое обучение на тысячах GPU или TPU --- стандартные инженерные решения.

Замечательное эмпирическое наблюдение состоит в том, что по мере роста модели и обучения на большем количестве данных она не просто запоминает обучающий корпус\,---\,она приобретает \textbf{эмерджентные способности}\index{emergent capabilities}, явно не присутствующие в обучающем сигнале. Модель, обученная только предсказывать следующее слово, учится:

\begin{itemize}[leftmargin=*]
  \item переводить между языками, которые она никогда не видела в парах;
  \item писать и отлаживать компьютерные программы;
  \item решать многошаговые математические задачи, показывая промежуточные вычисления;
  \item отвечать на фактические вопросы, используя знания, неявно закодированные в корпусе.
\end{itemize}

Wei et al.\ (2022) задокументировали, что некоторые способности появляются резко при определённых размерах моделей, что предполагает скачкообразные фазовые переходы в ландшафте потерь, а не плавную интерполяцию. Точный механизм эмерджентности остаётся активной областью исследований.

\subsection{Базовые модели и их ограничения}

Результатом предобучения является \textbf{базовая модель}\index{base model} (также называемая \textbf{фундаментальной моделью} или \textbf{предобученной моделью}). По заданному промпту базовая модель продолжает его в стиле того текста, с которым она научилась ассоциировать этот префикс. Такое поведение полезно для задач завершения текста, но проблематично при интерактивном использовании: на вопрос «Как испечь ржаной хлеб на закваске?» базовая модель может ответить большим количеством вопросов (как будто продолжая документ FAQ), а не полезным рецептом.

Что ещё важнее, базовая модель не имеет механизма отклонения вредных запросов, разграничения фактов и домыслов или принятия разговорного регистра, ожидаемого от ассистента. Эти ограничения мотивируют два последующих этапа обучения.

% ===========================================================
\section{Supervised Fine-Tuning}
\label{sec:sft}
\index{supervised fine-tuning (SFT)}

\textbf{Supervised fine-tuning} (SFT) адаптирует базовую модель для следования инструкциям, обучая её на тщательно отобранном датасете пар (инструкция, ответ). Процедура по существу идентична предобучению\,---\,предсказание следующего токена через кросс-энтропийную потерю\,---\,но на значительно меньшем, более качественном датасете демонстраций.

\subsection{Формат данных и шаблоны инструкций}

Данные SFT состоят из пар $(x^{(i)}, y^{(i)})$, где $x^{(i)}$ --- инструкция (запрос пользователя, вопрос или описание задачи), а $y^{(i)}$ --- высококачественный ответ, написанный или отобранный аннотатором-человеком или отфильтрованный из существующих высококачественных источников. На практике разговоры многоходовые, поэтому вход $x^{(i)}$ может включать полную историю диалога.

Для структурирования входа для модели \textbf{шаблон чата}\index{chat template} оборачивает каждый ход специальными токенами:
\begin{verbatim}
<|system|> Ты полезный ассистент. <|end|>
<|user|> Какова столица Франции? <|end|>
<|assistant|> Столицей Франции является Париж. <|end|>
\end{verbatim}
Модель обучается предсказывать только токены ассистента; токены пользователя и системы \emph{маскируются} при вычислении потерь, чтобы градиент не стимулировал модель имитировать речевые паттерны пользователя.

\subsection{Функция потерь SFT}

Потеря SFT --- та же кросс-энтропия, что~\eqref{eq:ntp-loss}, но применяемая только к токенам ответа:

\begin{equation}
\label{eq:sft-loss}
\mathcal{L}_{\mathrm{SFT}}(\theta) = -\frac{1}{N}\sum_{i=1}^{N}\sum_{t=1}^{T_i} \log p_\theta\!\left(y_t^{(i)} \mid x^{(i)}, y_{1:t-1}^{(i)}\right).
\end{equation}

Градиент~\eqref{eq:sft-loss} по $\theta$ вычисляется только по позициям, где $y_t$ --- токен ассистента. Эта асимметрия существенна: она учит модель генерировать хорошие ответы, не толкая её имитировать запросы пользователя или системные промпты.

\subsection{Датасеты и примеры}

\textbf{Stanford Alpaca} (Taori et al., 2023)\index{Alpaca} стал одним из первых широко воспроизведённых SFT-датасетов: 52~000 демонстраций следования инструкциям, сгенерированных с помощью промптинга \texttt{text-davinci-003} методом Self-Instruct (Wang et al., 2023). Несмотря на скромный размер, Alpaca показал, что 7-миллиардная модель LLaMA может приобрести полезное поведение следования инструкциям.

\textbf{T\"{u}lu} (Wang et al., 2023; Ivison et al., 2024)\index{Tülu} --- семейство моделей с настройкой по инструкциям, обученных на смеси отобранных датасетов с открытым исходным кодом. T\"{u}lu~3 (2024) использует более миллиона примеров следования инструкциям, собранных из разнообразных источников с тщательным дедублированием и фильтрацией качества, и включает как SFT, так и этапы дообучения на основе предпочтений. Он служит сильной базовой линией с открытым исходным кодом для изучения полного пайплайна постобучения.

\subsection{Смещение воздействия и потолок SFT}

SFT использует \textbf{принудительное обучение с учителем}\index{teacher forcing}: на каждом шаге $t$ во время обучения модель получает на вход \emph{истинный} предыдущий токен $y_{t-1}$, а не своё собственное предсказание $\hat{y}_{t-1}$. Во время вывода это предположение нарушается\,---\,модель должна обусловливаться на собственных (возможно, ошибочных) предыдущих выходах. Это несоответствие, известное как \textbf{смещение воздействия}\index{exposure bias} (Ranzato et al., 2015), может вызывать накопление ошибок: небольшая ошибка на шаге $t$ помещает модель в состояние, с которым она никогда не сталкивалась при обучении, делая дальнейшие ошибки более вероятными.

В более фундаментальном смысле SFT ограничен качеством своих демонстраций. Независимо от размера датасета SFT модель не может научиться производить ответы лучше, чем написавшие их аннотаторы. Для задач, где человеческие демонстрации дороги при масштабировании, или где пространство возможных хороших ответов богаче, чем может зафиксировать любой конечный датасет, этот потолок становится ограничивающим. Эти наблюдения мотивируют этап выравнивания\,---\,обучения на сравнениях предпочтений и сигналах вознаграждения, а не на фиксированных демонстрациях.

% ===========================================================
\section{Пайплайн постобучения}
\label{sec:post-training-pipeline}
\index{post-training pipeline}

Разработка современных языковых моделей следует трёхэтапному пайплайну, который можно понять как прогрессивный сдвиг обучающего сигнала от самообучающейся статистики к предпочтениям людей и верифицируемой корректности. Рисунок~\ref{fig:post-training-pipeline} иллюстрирует полный пайплайн.

\begin{figure}[ht]
\centering
\begin{tikzpicture}[
  node distance=0.7cm and 1.1cm,
  stage/.style={draw, rounded corners=5pt, minimum width=2.1cm, minimum height=1.0cm,
                text width=2.0cm, align=center, font=\small\bfseries},
  data/.style={draw, dashed, rounded corners=3pt, minimum width=1.9cm,
               minimum height=0.6cm, text width=1.8cm, align=center,
               font=\scriptsize, fill=gray!10},
  arr/.style={-{Stealth[length=5pt]}, thick},
  larr/.style={-{Stealth[length=4pt]}, dashed, gray},
  lbl/.style={font=\scriptsize\itshape, text width=2.0cm, align=center}
]

% ---- Верхний ряд: этапы пайплайна ----
\node[stage, fill=blue!12]   (pt)  {Пред-\\обучение};
\node[stage, fill=green!15,   right=0.6cm of pt]  (sft) {SFT};
\node[stage, fill=orange!18,  right=0.6cm of sft] (rm)  {Модель\\вознаг.};
\node[stage, fill=red!12,     right=0.6cm of rm]  (rl)  {RL-\\выравнивание};
\node[stage, fill=purple!12,  right=0.9cm of rl]  (fin) {Выровнен-\\ная LLM};

% ---- Стрелки между этапами ----
\draw[arr] (pt)  -- node[above, font=\scriptsize] {база}     (sft);
\draw[arr] (sft) -- node[above, font=\scriptsize] {SFT}      (rm);
\draw[arr] (rm)  -- node[above, font=\scriptsize] {$r_\psi$}  (rl);
\draw[arr] (rl)  --                                            (fin);

% ---- Метки потерь над этапами ----
\node[lbl, above=0.15cm of pt]  {NTP-потеря};
\node[lbl, above=0.15cm of sft] {маскир.~CE};
\node[lbl, above=0.15cm of rm]  {BT-потеря};
\node[lbl, above=0.15cm of rl]  {KL-штраф};

% ---- Узлы данных под этапами ----
\node[data, below=0.55cm of pt]  (dpt)  {Веб-корпус};
\node[data, below=0.55cm of sft] (dsft) {Демонстрации $(x,y)$};
\node[data, below=0.55cm of rm]  (drm)  {Предпочтения $(y_w,y_l)$};
\node[data, below=0.55cm of rl]  (drl)  {Роллауты + вознаграждение};

% ---- Стрелки данных ----
\draw[larr] (dpt)  -- (pt);
\draw[larr] (dsft) -- (sft);
\draw[larr] (drm)  -- (rm);
\draw[larr] (drl)  -- (rl);

\end{tikzpicture}
\caption{Полный пайплайн постобучения для большой языковой модели. Предобучение на веб-текстах создаёт базовую модель с широкими знаниями о мире. SFT на отобранных демонстрациях учит формату следования инструкциям. Моделирование вознаграждения обучает скалярному сигналу предпочтения по сравнениям людей. RL-выравнивание оптимизирует стратегию относительно этого вознаграждения, ограничивая дрейф штрафом KL. Главы~\ref{ch:practical-rlhf}--\ref{ch:open} подробно развивают каждый метод RL-выравнивания.}
\label{fig:post-training-pipeline}
\end{figure}

\subsection{Взаимодействие этапов}

Три этапа не являются независимыми: каждый существенно зависит от качества предыдущего.

\textbf{Предобучение устанавливает потолок знаний.} Этапы SFT и RL могут изменить поведение модели, но не могут привить знания, которых у базовой модели нет. Базовая модель, никогда не видевшая учебников по математике, не научится решать интегралы через настройку по инструкциям. Именно поэтому курирование данных предобучения столь значимо.

\textbf{SFT устанавливает поведенческий априор.} Модель SFT служит как отправной точкой, так и опорной стратегией $\pi_{\mathrm{ref}}$ для этапа RL. Слагаемое штрафа KL $\beta\,\KL{\pi_\theta}{\pi_{\mathrm{ref}}}$ (введённое в Главе~\ref{ch:reward-models}) привязывает стратегию RL вблизи распределения SFT. Плохо обученная модель SFT, следовательно, ограничивает эффективность RL-выравнивания.

\textbf{Моделирование вознаграждения определяет сигнал выравнивания.} Модель вознаграждения $r_\psi$ --- единственный источник обратной связи во время RL-обучения. Её точность, калибровка и охват распределения задач непосредственно определяют качество выровненной модели. Глава~\ref{ch:reward-models} подробно развивает теорию моделирования вознаграждения; Главы~\ref{ch:practical-rlhf}--\ref{ch:open} показывают, как модель вознаграждения используется в PPO, DPO и GRPO соответственно.

\subsection{Что оптимизирует каждый этап}

Краткое сравнение трёх этапов приведено в Таблице~\ref{tab:pipeline-comparison}.

\begin{table}[ht]
\centering
\caption{Сравнение трёх этапов постобучения.}
\label{tab:pipeline-comparison}
\footnotesize
\begin{tabular}{@{}p{1.9cm}p{2.9cm}p{2.4cm}p{1.9cm}p{2.8cm}@{}}
\toprule
\textbf{Этап} & \textbf{Данные} & \textbf{Цель} & \textbf{Парадигма} & \textbf{Ключевое ограничение} \\
\midrule
Предобучение & Неразмеченный текст & NLL следующего токена & Самообучение & Нет поведенческого контроля \\
SFT & Пары $(x, y)$ & NLL с маскировкой & Обучение с учителем & Ограничен качеством аннотаторов \\
RL-выравнивание & $(x, y_w, y_l)$ или $r(x,y)$ & KL-штрафованное вознаграждение & Обучение с подкреплением & Переоптимизация вознаграждения \\
\bottomrule
\end{tabular}
\end{table}

% ===========================================================
\section{Синтетические данные для постобучения}
\label{sec:synthetic-data}
\index{synthetic data}

Аннотирование людьми --- ограничивающий ресурс постобучения. Квалифицированный аннотатор производит около 20--50 высококачественных примеров следования инструкциям в час; один тренировочный прогон для передовой модели может потребовать миллионов. Этот арифметический разрыв побудил отрасль всё больше полагаться на \textbf{синтетические данные}: примеры, генерируемые моделями ИИ, а не непосредственно людьми-авторами. Этот раздел обозревает основные методы, их теоретические свойства и практические риски, которые они вносят.

\subsection{Почему синтетические данные необходимы}

Три структурных давления вынуждают использовать синтетические данные в масштабе.

\textbf{Объём.} Эффективное SFT и дообучение по предпочтениям требуют десятков тысяч до миллионов примеров, охватывающих разнообразные задачи, темы, уровни сложности и стили ответов. Аннотирование людьми в таком масштабе непомерно дорого: сборка одного миллиона высококачественных пар предпочтений по \$5 за сравнение обходится в \$5 миллионов, без учёта несогласия аннотаторов и накладных расходов на контроль качества.

\textbf{Охват.} Аннотаторы естественным образом кластеризуются вокруг знакомых областей. Редкие, но важные темы\,---\,нишевые научные субдисциплины, малоизвестные языки программирования, малоресурсные человеческие языки\,---\,систематически недопредставлены в сгенерированных людьми данных. Синтетическая генерация может быть направлена на эти пробелы с помощью инструкции.

\textbf{Калибровка сложности.} Обучение RL получает выгоду от примеров подходящего уровня сложности: задач, которые модель иногда, но не всегда решает, так чтобы градиент вознаграждения был информативным. Генерировать примеры с точно заданным целевым уровнем сложности практически невозможно с аннотаторами-людьми, но это можно сконструировать через синтетические пайплайны.

\subsection{Self-Instruct}
\index{Self-Instruct}

\textbf{Self-Instruct} (Wang et al., 2023) --- основополагающий метод синтетической генерации инструкций. Пайплайн начинается с небольшого начального набора из 175 написанных людьми инструкций, охватывающих разнообразный круг задач. На каждой итерации пайплайн запрашивает у большой языковой модели (\textbf{учителя}) несколько из этих исходных инструкций и просит сгенерировать новые инструкции, классифицировать их по типу задачи и построить пары «вход --- выход». Новые инструкции, слишком похожие на уже имеющиеся (по метрике ROUGE-L), отбрасываются; остальные добавляются в пул. Повторяя этот процесс, Wang et al.\ расширили исходный набор до 52~000 пар следования инструкциям\,---\,датасета Stanford Alpaca\,---\,при стоимости около \$500 в API-вызовах к \texttt{text-davinci-003}.

Статья Self-Instruct продемонстрировала, что 7-миллиардная модель LLaMA, дообученная на этих данных, сравнялась по способности следования инструкциям с моделью, обученной на частных написанных людьми данных, что говорит о том, что качество модели-учителя --- главное узкое место, а не само по себе различие между человеческим и синтетическим происхождением.

\subsection{Дистилляция из моделей-учителей}
\index{distillation!from teacher models}

Более прямой подход --- использовать более сильную модель (учителя) для генерации ответов, на которых затем обучается более слабая модель (ученик). Учитель может быть доступен через API (OpenAI, Anthropic, Google) или локально запущенную модель с открытыми весами. Процедура:

\begin{enumerate}[leftmargin=*]
  \item Собрать или сгенерировать набор промптов $\{x^{(i)}\}$, охватывающих желаемое распределение задач.
  \item Для каждого промпта сэмплировать один или несколько ответов от учителя: $y^{(i)} \sim \pi_{\mathrm{teacher}}(\cdot \mid x^{(i)})$.
  \item Дообучить ученика на результирующих парах $(x^{(i)}, y^{(i)})$ с помощью~\eqref{eq:sft-loss}.
\end{enumerate}

Дистилляция особенно эффективна для \emph{передачи формата}: обучения малой модели перенимать стиль, структуру и паттерны пошагового рассуждения большей модели. Серия Orca (Mukherjee et al., 2023) показала, что 13-миллиардная модель, обученная на подробных объяснительных ответах GPT-4, превзошла модели во много раз большего размера на определённых бенчмарках, что говорит о том, что ключевым ингредиентом является процесс рассуждения учителя, а не только его окончательный ответ.

Главное ограничение дистилляции --- правовое и коммерческое: большинство API передовых моделей запрещают использовать их выходы для обучения конкурирующих моделей. Учителя с открытыми весами (LLaMA~3, Mistral, Qwen) под это ограничение не подпадают и поэтому всё шире используются на практике.

\subsection{Evol-Instruct и нарастание сложности}
\index{Evol-Instruct}\index{WizardLM}

\textbf{Evol-Instruct} (Xu et al., 2023), введённый в рамках проекта \textbf{WizardLM}, заметил, что Self-Instruct склонен производить инструкции однородной (и часто низкой) сложности. Их идея --- \emph{эволюционировать} существующие инструкции, многократно применяя один из нескольких операторов мутации:

\begin{itemize}[leftmargin=*]
  \item \textbf{Добавить ограничения}: добавить требование, сужающее пространство ответов.
  \item \textbf{Углубить}: запросить более подробное или строгое рассмотрение.
  \item \textbf{Расширить}: обобщить инструкцию на более широкий контекст.
  \item \textbf{Добавить многошаговое рассуждение}: ввести требование, предполагающее поэтапный вывод.
  \item \textbf{Эволюция вширь}: сгенерировать совершенно новую, но связанную инструкцию.
\end{itemize}

Каждый шаг эволюции выполняется через промптинг модели-учителя. Результирующий датасет охватывает широкий спектр уровней сложности, что прямо выгодно для RL-обучения. WizardLM-2 (2024), обученный на данных Evol-Instruct в большом масштабе, достигает конкурентоспособного качества с моделями вдвое большего числа параметров на бенчмарках кодирования и математического рассуждения.

\subsection{Генерация на основе персон}
\index{persona-based generation}

Дополнительная стратегия увеличения разнообразия --- \textbf{генерация на основе персон}: вместо варьирования инструкции менять демографическую или профессиональную идентичность гипотетического пользователя. Модель-учитель запрашивается с описанием персоны, например «Вы аспирант второго года по вычислительной биологии» или «Вы старшеклассник, испытывающий трудности с алгеброй». Результирующие инструкции и ответы естественно кластеризуются вокруг разных словарей, уровней сложности и стилей общения.

Chan et al.\ (2024) показали, что генерация с учётом персон, применённая в масштабе с использованием 1 миллиарда уникальных синтетических персон, полученных из веб-текстов, даёт датасеты инструкций с существенно более высоким разнообразием, чем генерация без условия. Датасет T\"{u}lu~3 использует генерацию на основе персон как один из компонентов своей смеси данных.

\subsection{Синтез данных предпочтений}
\index{preference data synthesis}

Помимо демонстраций следования инструкциям, RL-выравнивание требует \textbf{данных предпочтений}: троек $(x, y_w, y_l)$, в которых человек (или модель) считает $y_w$ лучше $y_l$. Собирать такие сравнения от людей в масштабе ещё дороже, чем демонстрации, поскольку каждая аннотация требует от человека прочтения и сравнения двух полных ответов. Синтетическая генерация предпочтений решает это узкое место через три взаимодополняющих механизма.

\textbf{Множественное сэмплирование ответов и оценка LLM.} Наиболее распространённый подход:
\begin{enumerate}[leftmargin=*]
  \item Для каждого промпта $x$ сэмплировать $k \geq 2$ ответов от текущей или модели-учителя.
  \item Оценить каждый ответ с помощью \textbf{модели-судьи}\index{LLM-as-judge} (обычно GPT-4 или судьи с открытыми весами) по рубрике.
  \item Формировать пары предпочтений, принимая $y_w$ как ответ с наибольшим баллом и $y_l$ как любой ответ с меньшим баллом (или специально с наименьшим).
\end{enumerate}

\textbf{UltraFeedback} (Cui et al., 2023)\index{UltraFeedback} --- канонический датасет, собранный этим методом: 64~000 промптов, каждый связан с четырьмя ответами от различных моделей, оцененными GPT-4 по четырём рубрикам (следование инструкциям, правдивость, честность, полезность). Результирующий датасет предпочтений использовался для обучения десятков выровненных моделей, включая серию Zephyr.

\textbf{Constitutional AI и RLAIF.} Bai et al.\ (2022) из Anthropic представили \textbf{Constitutional AI} (CAI)\index{Constitutional AI (CAI)}, в котором сама модель ИИ критикует и пересматривает свои ответы в соответствии с фиксированным набором принципов («конституцией»). Цикл критика-пересмотр генерирует улучшенные ответы без вмешательства людей, а результирующие пары (исходный, пересмотренный) можно использовать как данные предпочтений. Тесно связанный \textbf{RLAIF} (Reinforcement Learning from AI Feedback; Lee et al., 2023)\index{RLAIF} использует разметчик ИИ вместо разметчиков-людей для аннотирования того, какой из двух ответов предпочтительнее. Lee et al.\ показали, что RLAIF и человеческий RLHF дают модели сопоставимого качества на нескольких стандартных бенчмарках, что говорит о том, что узким местом в выравнивании является не источник меток предпочтений (человек или ИИ), а их точность и охват.

\textbf{Оценка по рубрикам.} Вместо единственного общего предпочтения современные пайплайны часто просят судью оценить по нескольким измерениям\,---\,корректность, лаконичность, безопасность, соответствие инструкции\,---\,и агрегировать. Рубрики также могут быть задачеспецифичными: для математики рубрика проверяет, верен ли окончательный ответ и валидны ли промежуточные шаги. T\"{u}lu~3 использует рубрики с верификаторами для математических и кодовых задач, полностью заменяя модель-судью в областях, где корректность проверяема.

\subsection{Фильтрация данных и контроль качества}
\index{data quality filtering}

Синтетическая генерация данных дешева, но шумна. Механизмы контроля качества существенны для предотвращения деградации производительности модели от низкокачественных примеров. Стандартные техники включают:

\begin{itemize}[leftmargin=*]
  \item \textbf{Фильтрация по перплексии}: отклонять ответы, которым текущая модель присваивает чрезвычайно высокую или низкую перплексию (первое указывает на деформированный текст; второе --- на почти запомненный шаблон).
  \item \textbf{Фильтрация по длине}: отбрасывать ответы за пределами разумного диапазона длин (очень короткие часто лишены деталей; очень длинные часто заполняются пустыми словами).
  \item \textbf{Дедупликация}: удалять почти-дубликаты с помощью MinHash или косинусного сходства на основе эмбеддингов для предотвращения запоминания модели повторяющихся паттернов.
  \item \textbf{Фильтрация по модели вознаграждения}: сохранять только ответы, которым модель вознаграждения присваивает балл выше порога. Это создаёт свойство самосогласованности: данные, используемые для обучения, уже считаются хорошими моделью вознаграждения.
  \item \textbf{Фильтрация по сложности}: для RL-обучения сохранять только примеры в «обучаемой зоне»\,---\,где частота успеха модели составляет от 20\% до 80\%. Задачи, которые модель всегда решает, не дают градиента; задачи, которые она никогда не решает, вносят только шум.
\end{itemize}

\subsection{Коллапс модели и меры противодействия}
\index{model collapse}

Фундаментальная проблема обучения на синтетических данных --- \textbf{коллапс модели}\index{model collapse} (Shumailov et al., 2024): если модели обучаются на своих собственных выходах (или выходах аналогично обученных моделей), обучающее распределение может прогрессивно сужаться, заставляя модель терять широту и разнообразие, присутствующие в исходном написанном людьми корпусе. Математически, если обучающее распределение в поколении $g$ --- $p_g$, а распределение синтетических данных является приближением $p_{g-1}$, итерирование этого процесса вызывает геометрическое уменьшение дисперсии $p_g$, в конечном счёте коллапсируя к узкой моде.

Определены четыре меры противодействия:

\begin{enumerate}[leftmargin=*]
  \item \textbf{Смешивание с человеческими данными.} Всегда сохранять долю подлинного написанного людьми текста в обучающей смеси. Доля не обязана быть большой: Gerstgrasser et al.\ (2024) показали, что даже 5--10\% человеческих данных достаточно для теоретического предотвращения коллапса.
  \item \textbf{Разнообразный ансамбль учителей.} Использование нескольких моделей-учителей с различными архитектурами и историями обучения вносит дистрибуционное разнообразие, которое единственный учитель не может предоставить.
  \item \textbf{Агрессивная дедупликация.} Почти-дублированные синтетические примеры --- механизм распространения коллапса моды; их удаление разрывает положительную обратную связь.
  \item \textbf{Водяные знаки и отслеживание происхождения.} Маркировка синтетических данных при генерации позволяет последующим кураторам мониторировать долю синтетических данных и ограничивать её по мере необходимости.
\end{enumerate}

\subsection{Ключевые синтетические датасеты}
\index{datasets!synthetic}

Таблица~\ref{tab:synthetic-datasets} суммирует наиболее влиятельные синтетические датасеты в литературе по постобучению.

\begin{table}[ht]
\centering
\caption{Основные синтетические датасеты, используемые в постобучении.}
\label{tab:synthetic-datasets}
\begin{tabular}{lllll}
\toprule
\textbf{Датасет} & \textbf{Год} & \textbf{Размер} & \textbf{Тип} & \textbf{Метод} \\
\midrule
Stanford Alpaca     & 2023 & 52К    & SFT          & Self-Instruct (GPT-3.5) \\
UltraFeedback       & 2023 & 256К пар & Предпочтения  & LLM-судья (GPT-4) \\
T\"{u}lu~3          & 2024 & 1М+    & SFT + Предп.  & Смесь, персоны, рубрики \\
OpenThoughts~3      & 2025 & 1М+    & SFT (CoT)    & Дистилляция (DeepSeek-R1) \\
\bottomrule
\end{tabular}
\end{table}

\textbf{OpenThoughts~3}\index{OpenThoughts} (2025) воплощает нарождающийся класс \emph{ориентированных на рассуждение} синтетических датасетов: один миллион цепочек рассуждений для математических, научных и кодовых задач, дистиллированных из DeepSeek-R1. Обучение базовой модели на OpenThoughts~3 даёт модель, демонстрирующую многошаговое рассуждение в формате $\langle\text{thinking}\rangle$, существенно улучшая производительность на олимпиадной математике (AIME, MATH) и бенчмарках кодирования (HumanEval).

% ===========================================================
\section{Оценка языковых моделей}
\label{sec:lm-eval}
\index{evaluation!language models}

Оценка языковых моделей существенно сложнее, чем оценка классических RL-агентов. Стратегия в мире клеток либо достигает цели, либо нет; качество разговорного ответа не поддаётся такой бинарной характеристике. В этом разделе обозреваются основные парадигмы оценки.

\subsection{Автоматизированные бенчмарки}
\index{benchmarks}

Автоматизированные бенчмарки обеспечивают дешёвую, воспроизводимую и объективную оценку. Наиболее широко используемые включают:

\begin{description}[leftmargin=!, labelwidth=3.5cm]
  \item[MMLU] (Hendrycks et al., 2021) Massive Multitask Language Understanding: 57 академических предметов от уровня старшей школы до профессионального, в формате множественного выбора. Проверяет широту фактических знаний.
  \item[HellaSwag] (Zellers et al., 2019) Логический вывод на основе здравого смысла: выбрать наиболее вероятное продолжение описания деятельности. Проверяет заземлённое рассуждение на основе здравого смысла.
  \item[GSM8K] (Cobbe et al., 2021) Математика уровня начальной школы: 8500 многошаговых арифметических текстовых задач. Ответы верифицируются проверкой окончательного числового результата.
  \item[HumanEval] (Chen et al., 2021) Генерация кода: 164 задачи на программирование на Python, оцениваемые запуском юнит-тестов против сгенерированного кода.
  \item[MATH] (Hendrycks et al., 2021) Олимпиадная математика: 12~500 задач уровня AMC/AIME, требующих символьных манипуляций и многошагового доказательства.
\end{description}

Автоматизированные бенчмарки имеют критический режим отказа: модели могут быть \textbf{загрязнены}\index{benchmark contamination}, если данные оценки появляются в корпусе предобучения. Оценка GPT-4 на GSM8K улучшилась на 7 процентных пунктов при реконструкции оценочного разбиения с новыми формулировками задач, что говорит о частичном запоминании. Ответственная оценка требует проверки n-граммового пересечения между бенчмарком и обучающими данными.

\subsection{Оценка людьми}
\index{human evaluation}

Для открытых задач, где не существует истинного ответа\,---\,творческое письмо, диалог в открытом домене, реферирование\,---\,оценка людьми остаётся золотым стандартом. Аннотаторов обычно просят сравнить два ответа бок о бок (A/B-оценка) и указать, какой лучше по одному или нескольким измерениям. Платформы, такие как Chatbot Arena (Chiang et al., 2024), собирают попарные предпочтения в масштабе от пользователей-добровольцев, производя Эло-стиль ранжирования моделей.

Оценка людьми дорога, медленна и страдает систематическими смещениями: аннотаторы предпочитают более длинные ответы (смещение по многословности), ответы, соглашающиеся с их предварительными взглядами (подхалимство), и ответы от моделей, которые они считают более способными (якорение). Контроль за этими смещениями требует тщательного экспериментального дизайна и больших выборок.

\subsection{LLM как судья}
\index{LLM-as-judge}

\textbf{LLM-as-judge} (Zheng et al., 2023) использует сильную языковую модель (обычно GPT-4) как автоматизированного оценщика, заменяя аннотаторов-людей в A/B-сравнениях. Судье подают промпт с инструкцией, двумя ответами и структурированной рубрикой, после чего он должен вынести суждение о предпочтении и кратко его объяснить. MT-Bench и AlpacaEval 2.0 используют эту парадигму для ранжирования моделей, настроенных на следование инструкциям.

LLM-as-judge наследует смещения модели-судьи и вносит \textbf{смещение самоулучшения}: модели от того же разработчика, что и судья, как правило, оцениваются более благоприятно. Меры противодействия включают использование нескольких судей, запрос объяснений в виде цепочки рассуждений перед окончательным вердиктом и перестановку позиций (оценку каждой пары в обоих порядках с усреднением).

\subsection{Оценка модели вознаграждения}
\index{reward model!evaluation}

Когда конечным использованием оценки является выбор между кандидатными моделями во время RL-обучения, обученная модель вознаграждения $r_\psi$ используется как оценщик. Точность модели вознаграждения измеряется как доля отложенных пар предпочтений, для которых она правильно предсказывает предпочтительный с точки зрения человека ответ. Как отмечено в Главе~\ref{ch:reward-models}, хорошо откалиброванные модели вознаграждения достигают 70--80\% точности на отложенных человеческих предпочтениях. Помимо точности, распределение оценок вознаграждения и их корреляция с независимыми человеческими рейтингами мониторируются для обнаружения систематических смещений.

% ===========================================================
\section{Резюме}
\label{sec:lm-summary}

В этой главе была представлена инфраструктура языкового моделирования, лежащая в основе методов RLHF, развиваемых в последующих главах.

Языковая модель --- это авторегрессионное распределение вероятностей на последовательностях токенов, реализованное с использованием архитектуры Трансформера. Предобучение на веб-текстах с использованием предсказания следующего токена прививает широкие мировые знания и эмерджентные способности; supervised fine-tuning на отобранных парах инструкция-ответ учит поведенческому формату, ожидаемому от ассистента. Затем этап RL-выравнивания со штрафом KL оптимизирует сигнал вознаграждения\,---\,выученный из сравнений предпочтений людей или верифицируемых исходов задач\,---\,для поднятия модели выше потолка, установленного демонстрациями.

Ключевым концептуальным мостиком является MDP уровня токенов (Определение~\ref{def:token-mdp}): генерация языка точно отображается на MDP-фреймворк Глав~\ref{ch:mdp}--\ref{ch:actor-critic}, где языковая модель --- стратегия, словарь --- пространство действий, а терминальное вознаграждение от человека или модели вознаграждения. Это соответствие оправдывает прямое применение методов градиента стратегии к обучению языковых моделей.

Синтетические данные стали незаменимым компонентом каждого этапа постобучения: Self-Instruct и Evol-Instruct для разнообразия инструкций SFT, пайплайны LLM-судьи для данных предпочтений в масштабе и дистилляция рассуждений для дообучения по цепочке рассуждений. Коллапс модели --- реальный риск, но его можно смягчить, смешивая человеческие данные, используя разнообразные ансамбли учителей и агрессивную дедупликацию.

Последующие главы подробно развивают три основных алгоритма RL-выравнивания. Глава~\ref{ch:practical-rlhf} рассматривает RLHF на основе PPO, пайплайн, ставший пионером InstructGPT и используемый в GPT-4 и Claude~2. Последующие главы охватывают DPO, устраняющий явную модель вознаграждения, и GRPO, снижающий дисперсию через вычитание групповой относительной базовой линии и центральный для моделей рассуждений DeepSeek. В каждом случае MDP уровня токенов, установленный здесь, предоставляет общий формальный язык.

% ===========================================================
\section*{Упражнения}
\addcontentsline{toc}{section}{Упражнения}

\begin{exercise}
\label{ex:lm-perplexity}
Языковая модель присваивает следующие вероятности токенам предложения «The cat sat»: $p(\text{The}) = 0.05$, $p(\text{cat} \mid \text{The}) = 0.02$, $p(\text{sat} \mid \text{The cat}) = 0.08$.
\begin{enumerate}
  \item Вычислите кросс-энтропийную потерю~\eqref{eq:ntp-loss} для этого предложения.
  \item Вычислите перплексию~\eqref{eq:perplexity}.
  \item Если вместо этого используется равномерное распределение на словаре размером $|\mathcal{V}| = 50~000$, какова перплексия? Сравните и интерпретируйте.
\end{enumerate}
\end{exercise}

\begin{exercise}
\label{ex:token-mdp-size}
Рассмотрим MDP уровня токенов (Определение~\ref{def:token-mdp}) с размером словаря $|\mathcal{V}| = 100~000$ и максимальной длиной ответа $T = 512$ токенов.
\begin{enumerate}
  \item Дайте верхнюю оценку $|\cS|$, числа возможных состояний. Сколько бит памяти потребовалось бы для хранения таблицы $Q(s, a)$ для всех пар $(s, a)$, если каждая запись занимает 4 байта?
  \item Объясните, почему методы на основе ценности, такие как DQN, неосуществимы для этого MDP.
  \item Теорема о градиенте стратегии (Глава~\ref{ch:pg}) требует только $\nabla_\theta \log \pi_\theta(a \mid s)$, а не Q-таблицы. Почему это делает методы на основе градиента стратегии трактуемыми несмотря на огромное пространство состояний-действий?
\end{enumerate}
\end{exercise}

\begin{exercise}
\label{ex:causal-mask}
Причинная маска в Трансформере только-декодера гарантирует, что позиция $i$ не может обращать внимание на позиции $j > i$.
\begin{enumerate}
  \item Запишите маскированный вес внимания $\alpha_{ij}$ формулой, используя индикаторную функцию $\ind[j \leq i]$.
  \item Объясните, почему эта маска эквивалентна марковскому свойству MDP уровня токенов: предсказание $y_t$ зависит только от $(x, y_{1:t-1})$, но не от любого будущего токена.
  \item Была бы необходима причинная маска для языковой модели, используемой только как \emph{классификатор} (предсказывающий метку для всей последовательности)? Обоснуйте ответ.
\end{enumerate}
\end{exercise}

\begin{exercise}
\label{ex:sft-exposure-bias}
Проблема exposure bias возникает потому, что SFT использует принуждение учителем во время обучения, но во время вывода опирается на собственные выходы модели.
\begin{enumerate}
  \item Формализуйте несоответствие: пусть $p_{\mathrm{teacher}}$ --- распределение обучения с принуждением учителя по контекстам на шаге $t$, а $p_{\mathrm{model}}$ --- соответствующее распределение вывода. Покажите, что они совпадают тогда и только тогда, когда модель идеальна (присваивает вероятность~1 верному следующему токену на каждом шаге).
  \item Запланированное сэмплирование (Bengio et al., 2015) заменяет истинный токен собственным предсказанием модели с вероятностью $\epsilon_t$, где $\epsilon_t$ возрастает от 0 до 1 в ходе обучения. Опишите одно преимущество и один недостаток этого подхода по сравнению со стандартным принуждением учителем.
  \item Объясните, почему RL-дообучение (которое обучается на собственных роллаутах модели) не страдает от смещения воздействия.
\end{enumerate}
\end{exercise}

\begin{exercise}
\label{ex:chinchilla}
Закон масштабирования Chinchilla утверждает, что оптимальное по вычислениям обучение требует $D \approx 20N$ токенов для модели с $N$ параметрами. Предположим вычислительный бюджет $C = 6 \times 10^{23}$ FLOPs и что один обучающий FLOP стоит приблизительно $C \approx 6ND$ FLOPs.
\begin{enumerate}
  \item Выведите оптимальный по вычислениям размер модели $N^*$ и количество токенов $D^*$ как функции $C$.
  \item С приведённым выше бюджетом вычислите $N^*$ и $D^*$ численно. Сравните с публично известной конфигурацией LLaMA~3 8B (8B параметров, 15T токенов). Является ли LLaMA~3 8B оптимальной по вычислениям согласно Chinchilla?
  \item Анализ Chinchilla минимизирует потерю при обучении. Аргументируйте, почему практик, оптимизирующий \emph{стоимость вывода}, может рационально выбрать меньшую, более переобученную модель.
\end{enumerate}
\end{exercise}

\begin{exercise}
\label{ex:self-instruct}
Пайплайн Self-Instruct (Wang et al., 2023) засеивает генерацию инструкций 175 написанными людьми примерами и развивает датасет через промптинг модели-учителя.
\begin{enumerate}
  \item Опишите один механизм, с помощью которого качество начального набора влияет на качество сгенерированного датасета.
  \item Сходство ROUGE-L используется для фильтрации почти-дублированных инструкций. Приведите пример двух инструкций, которые ROUGE-L пометит как похожие, но которые подлинно различны в своих когнитивных требованиях, и наоборот (две инструкции, которые ROUGE-L оценивает как несхожие, но которые семантически почти идентичны).
  \item Предположим, что модель-учитель, используемая в Self-Instruct, систематически склонна к определённым стилям ответа (например, всегда использует маркированные списки). Как это смещение распространится на студенческую модель, обученную на результирующих данных?
\end{enumerate}
\end{exercise}

\begin{exercise}
\label{ex:model-collapse}
Коллапс модели означает прогрессивное сужение выходного распределения модели при итеративном обучении на её собственных синтетических выходах.
\begin{enumerate}
  \item Пусть $p_0$ --- исходное распределение человеческих данных, $\hat{p}_g$ --- распределение, выученное в поколении $g$, аппроксимированное обучением на выборках из $\hat{p}_{g-1}$. Предположим, что каждый шаг обучения вносит ошибку аппроксимации, ограниченную $\epsilon$ по расстоянию суммарной вариации: $\mathrm{TV}(\hat{p}_g, p_{g-1}) \leq \epsilon$. Дайте верхнюю оценку $\mathrm{TV}(\hat{p}_G, p_0)$ после $G$ поколений.
  \item Shumailov et al.\ (2024) показали, что хвосты распределения (редкий, но важный контент) теряются первыми. Объясните интуитивно, почему это так, используя интерпретацию через дисперсию при коллапсе модели.
  \item Предложите стратегию смешивания данных, которая теоретически предотвращает неограниченный рост оценки суммарной вариации при $G \to \infty$, предполагая, что доля $\lambda$ обучающих данных в каждом поколении --- подлинные данные людей.
\end{enumerate}
\end{exercise}

\begin{exercise}
\label{ex:llm-as-judge}
LLM-as-judge использует способную модель для оценки ответов, заменяя аннотаторов-людей.
\begin{enumerate}
  \item Опишите \emph{смещение по многословности} в системах LLM-as-judge и предложите меры противодействия посредством инженерии промптов.
  \item Позиционное смещение относится к тенденции судей-LLM предпочитать ответ, представленный первым (или вторым). Разработайте протокол оценки, устраняющий позиционное смещение, и объясните, почему ваш протокол несмещён.
  \item Предположим, что модель вознаграждения, обученная на предпочтениях, сгенерированных судьёй-LLM, впоследствии используется в PPO для производства выровненной модели. Проследите, как систематическое смещение в модели-судье распространяется по пайплайну, влияя на поведение итоговой выровненной модели.
\end{enumerate}
\end{exercise}

\begin{exercise}
\label{ex:constitutional-ai}
Constitutional AI (Bai et al., 2022) использует набор принципов для руководства самокритикой и пересмотром.
\begin{enumerate}
  \item Запишите цикл критика-пересмотр как двухшаговую процедуру. Какова роль конституции? Чем это отличается от стандартного SFT?
  \item RLAIF использует тот же принцип, но на этапе разметки предпочтений: модель просят оценить, какой из двух ответов лучше соответствует конституции. Покажите, что если модель-судья идеальна (всегда выбирает подлинно лучший ответ), то RLAIF и человеческий RLHF дают одну и ту же оптимальную стратегию по модели Bradley--Terry (Глава~\ref{ch:reward-models}).
  \item Каков основной режим отказа RLAIF при несовершенной модели-судье, и как он связан с переоптимизацией вознаграждения?
\end{enumerate}
\end{exercise}

\begin{exercise}
\label{ex:mdp-mapping}
Рассмотрим простую задачу генерации языка: модель должна произвести бинарную строку длины $T = 3$ (например, «010»), где вознаграждение равно~$+1$, если число единиц равно числу нулей, и $-1$ иначе.
\begin{enumerate}
  \item Запишите полный MDP уровня токенов: перечислите все состояния, действия, переходы и вознаграждения.
  \item Вычислите оптимальную стратегию $\pi^*$ с помощью динамического программирования (задача достаточно мала для ручного вычисления).
  \item Для $T = 3$ существует $2^3 = 8$ возможных ответов. Покажите, что цель RLHF~\eqref{eq:rlhf-objective} равна $\Prob_{\pi_\theta}(\mbox{сбалансированная строка}) - \Prob_{\pi_\theta}(\mbox{несбалансированная строка})$, и вычислите её значение при равномерной стратегии $\pi_\theta(0 \mid s) = \pi_\theta(1 \mid s) = 0.5$ для всех $s$.
\end{enumerate}
\end{exercise}
